% Part: computability
% Chapter: computability-theory
% Section: equiv-ce-defs

\documentclass[../../../include/open-logic-section]{subfiles}

\begin{document}

\olfileid{cmp}{thy}{eqc}

\olsection[संगणनक्षमपणे प्रगणनीय संचांच्या व्याख्या]{संगणनक्षमपणे
  प्रगणनीय संचांच्या समतुल्य व्याख्या}


संगणनक्षमपणे प्रगणनीय असण्याचा अर्थ सांगणारी पुढील काही महत्त्वाची
समतुल्य विधाने आहेत.

\begin{thm}
\ollabel{thm:ce-equiv}
$S$ हा नैसर्गिक संख्यांचा संच असू द्या. मग पुढील विधाने समतुल्य आहेत:
\begin{enumerate}
\item\ollabel{case:ce} $S$ संगणनक्षमपणे प्रगणनीय आहे.
\item\ollabel{case:ran-pc} $S$ ही एखाद्या \emph{आंशिक}
  संगणनक्षम फलनाची व्याप्ती आहे.
\item\ollabel{case:ran-prim} $S$ रिकामा आहे किंवा एखाद्या आदिम
  पुनरावर्ती फलनाची व्याप्ती आहे.
\item\ollabel{case:ce-domain} $S$ हे एखाद्या आंशिक संगणनक्षम
  फलनाचे \emph{परिभाषाक्षेत्र} आहे.
\end{enumerate}
\end{thm}

\begin{explain}
पहिली तीन विधाने सांगतात की कोणत्याही रिकाम्या नसलेल्या
संगणनक्षमपणे प्रगणनीय संचाचे प्रगणन संगणनक्षम फलनाने, आंशिक
संगणनक्षम फलनाने किंवा आदिम पुनरावर्ती फलनाने करता येते.
चौथे विधान सांगते की $S$ संगणनक्षमपणे प्रगणनीय असेल, तर
एखाद्या निर्देशांक~$e$ साठी
\[
S = \Setabs{x}{\cfind{e}(x) \fdefined}.
\]
दुसऱ्या शब्दांत, $S$ हा अशा आदानांचा संच आहे ज्यांवर
$\cfind{e}$ चे संगणन थांबते. म्हणून संगणनक्षमपणे प्रगणनीय
संचांना कधी कधी \emph{अर्धनिर्णेय} म्हणतात: संख्या त्या संचात
असेल, तर शेवटी ``होय'' असे उत्तर मिळते; पण ती संचात नसेल,
तर ``नाही'' असे उत्तर कधीच मिळत नाही!{}
\end{explain}

\begin{proof}
प्रत्येक आदिम पुनरावर्ती फलन संगणनक्षम असते आणि प्रत्येक
संगणनक्षम फलन आंशिक संगणनक्षमही असते. म्हणून
\olref{case:ran-prim} वरून \olref{case:ce} मिळते, आणि
\olref{case:ce} वरून~\olref{case:ran-pc} मिळते.
($S$ रिकामा असेल, तर $S$ कुठेही परिभाषित नसलेल्या आंशिक
संगणनक्षम फलनाची व्याप्ती असतो, हे लक्षात घ्या.)
\olref{case:ran-pc} वरून \olref{case:ran-prim} मिळते हे
दाखवले, तर पहिली तीन विधाने समतुल्य ठरतील.

म्हणून $S$ ही आंशिक संगणनक्षम फलन $\cfind{e}$ ची व्याप्ती आहे
असे समजा. $S$ रिकामा असेल, तर काम झाले. अन्यथा $a$ हा~$S$
मधील कोणताही घटक घ्या. क्लिनी यांच्या प्रमाण रूप प्रमेयावरून
पुढीलप्रमाणे लिहिता येते:
\[
\cfind{e}(x) = U(\umin{s}{T(e, x, s)}).
\]
विशेषतः, $\cfind{e}(x) \fdefined$ असून त्याचे मूल्य $= y$
असेल तेव्हाच असा $s$ असतो की $T(e, x, s)$ आणि $U(s) = y$.
$f(z)$ ची व्याख्या पुढीलप्रमाणे करा:
\[
f(z) = \begin{cases}
  U((z)_1) & \text{जर $T(e, (z)_0, (z)_1)$ असेल तर} \\
  a        & \text{अन्यथा.}
\end{cases}
\]
मग $f$ आदिम पुनरावर्ती आहे, कारण $T$ आणि $U$
तशी आहेत. \iftag{TMs}{ट्यूरिंग यंत्रांच्या भाषेत, $z$ हे
  $\tuple{(z)_0, (z)_1}$ या जोडीचा संकेतांक असेल आणि
  $(z)_1$ हे यंत्र~$M_e$ चे आदान $(z)_0$ वरील थांबणारे
  संगणन असेल, तर $f$ त्या संगणनाचे निर्गत परत करते;
  अन्यथा ते~$a$ परत करते.}

$S$ ही~$f$ ची व्याप्ती आहे हे दाखवायचे आहे; म्हणजे प्रत्येक
नैसर्गिक संख्या~$y$ साठी $y \in S$ असेल तेव्हाच ती~$f$
च्या व्याप्तीत असेल. प्रथम $y \in S$ असे समजा. मग $y$ ही
$\cfind{e}$ च्या व्याप्तीत आहे. म्हणून काही $x$ आणि~$s$
साठी $T(e,x,s)$ लागू होतो आणि $U(s) = y$. त्यामुळे
$y = f(\tuple{x,s})$. उलट, $y$ ही~$f$ च्या व्याप्तीत
असेल, तर एकतर $y = a$ असेल, किंवा एखाद्या~$z$ साठी
$T(e,(z)_0,(z)_1)$ लागू होईल आणि $U((z)_1) = y$ असेल.
दुसऱ्या प्रकरणात $\cfind{e}((z)_0) \fdefined =
y$ असल्याने, दोन्ही प्रकारे $y$ ही~$S$ मध्येच असते.
%READERNOTE{OLCMP-017 — गोठवलेल्या इंग्रजीतील व्याप्तीच्या उलट दिशेच्या युक्तिवादात x अनिर्दिष्ट आहे. येथे साक्षीदार z चा पहिला घटक (z)_0 हा योग्य आदान आहे; मराठीत तो स्पष्ट केला आहे.}

($\cfind{e}(x) \fdefined = y$ या चिन्हांकनाचा अर्थ
``$\cfind{e}(x)$ परिभाषित आहे आणि $y$ च्या समान आहे'' असा आहे.
त्याऐवजी $\cfind{e}(x) =
y$ असेही लिहिता येईल; पण आपण आंशिक फलनाशी व्यवहार करत आहोत,
याची आठवण करून देण्यासाठी अतिरिक्त बाण कधी कधी उपयुक्त असतो.)

\olref{thm:ce-equiv} ची सिद्धता पूर्ण करण्यासाठी
\olref{case:ce} आणि~\olref{case:ce-domain} समतुल्य आहेत
हे दाखवणे पुरेसे आहे. प्रथम \olref{case:ce} वरून
\olref{case:ce-domain} मिळते हे पाहू. संच रिकामा असेल,
तर कुठेही परिभाषित नसलेल्या आंशिक संगणनक्षम फलनाचे तो
परिभाषाक्षेत्र आहे. आता तो रिकामा नाही असे समजा. मग
$S$ ही संगणनक्षम फलन~$f$ ची व्याप्ती आहे, म्हणजे
\[
S = \Setabs{y}{\text{अशा काही $x$ साठी, } f(x) = y}.
\]
पुढीलप्रमाणे घ्या:
\[
g(y) = \umin{x}{(f(x) = y)}.
\]
मग $g$ आंशिक संगणनक्षम आहे आणि काही~$x$ साठी $f(x) = y$
असेल तेव्हाच $g(y)$ परिभाषित असते. दुसऱ्या शब्दांत, $g$
चे परिभाषाक्षेत्र ही~$f$ ची व्याप्ती आहे.
\iftag{TMs}{ट्यूरिंग यंत्रांच्या भाषेत: $S$ चे घटक प्रगणन
करणारे ट्यूरिंग यंत्र~$F$ दिले असेल, तर $G$ हे ट्यूरिंग यंत्र
दिलेला घटक संचात आहे का हे तपासण्यासाठी~$F$ च्या निर्गतांमध्ये
शोध घेते. तो घटक सापडल्यास ते थांबते; न सापडल्यास
ते कायम शोधत राहते. अशा प्रकारे ते~$S$ अर्धनिर्णीत करते.}{}
%READERNOTE{OLCMP-018 — गोठवलेल्या इंग्रजीत या दिशेची सिद्धता थेट S ही संगणनक्षम फलनाची व्याप्ती आहे असे धरते; पण व्याख्येनुसार रिकामा S देखील c.e. आहे आणि तो सर्वत्र परिभाषित फलनाची व्याप्ती असू शकत नाही. रिकाम्या संचासाठी कुठेही परिभाषित नसलेल्या फलनाचा स्वतंत्र आधार मराठीत दिला आहे.}

शेवटी \olref{case:ce-domain} वरून~\olref{case:ce} मिळते
हे दाखवू. $S$ हे आंशिक संगणनक्षम फलन~$\cfind{e}$ चे
परिभाषाक्षेत्र आहे असे समजा, म्हणजे
\[
S = \Setabs{x}{\cfind{e}(x) \fdefined}.
\]
$S$ रिकामा असेल, तर काम झाले; अन्यथा $a$ हा~$S$
मधील कोणताही घटक घ्या. $f$ ची पुढीलप्रमाणे व्याख्या करा:
\[
f(z) = \begin{cases}
(z)_0 & \text{जर $T(e,(z)_0,(z)_1)$ असेल तर} \\
a & \text{अन्यथा.}
\end{cases}
\]
मग वरच्याप्रमाणे, संख्या $x$ ही~$f$ च्या व्याप्तीत असेल
तेव्हाच $\cfind{e}(x) \fdefined$, म्हणजे $x \in S$.
\iftag{TMs}{ट्यूरिंग यंत्रांच्या भाषेत: $S$ अर्धनिर्णीत
करणारे यंत्र $M_e$ दिले असेल, तर सर्व शक्य ट्यूरिंग-यंत्र
संगणनांमधून क्रमाने जाऊन थांबणाऱ्या संगणनांना अनुरूप आदाने
परत करत~$S$ चे प्रगणन करा.}{}
\end{proof}

\olref{thm:ce-equiv} मधील विधान~\olref{case:ce-domain}
संगणनक्षमपणे प्रगणनीय संचांचे प्रगणन करण्याचा सोयीचा मार्ग देते:
प्रत्येक~$e$ साठी $W_e$ हे $\cfind{e}$ चे परिभाषाक्षेत्र दर्शवू,
म्हणजे
\[
W_e = \Setabs{x}{\cfind{e}(x) \fdefined}.
\]
मग $A$ हा कोणताही संगणनक्षमपणे प्रगणनीय संच असेल, तर
एखाद्या~$e$ साठी $A = W_e$.

पुढील प्रमेय संगणनक्षमपणे प्रगणनीय संचांचे आणखी एक वैशिष्ट्य
सांगते.

\begin{thm}
\ollabel{thm:exists-char}
$S$ हा संच संगणनक्षमपणे प्रगणनीय असेल तेव्हाच असा
संगणनक्षम संबंध $R(x,y)$ असतो की
\[
S = \Setabs{ x }{ \lexists[y][R(x,y)] }.
\]
\end{thm}

\begin{proof}
प्रथम $S$ संगणनक्षमपणे प्रगणनीय आहे असे समजा. मग काही
$e$ साठी $S = W_e$. या~$e$ साठी $S$ असे लिहिता येते:
\[
S = \Setabs{ x }{ \lexists[y][T(e, x, y)] }.
\]
उलट, $S = \Setabs{ x }{
  \lexists[y][R(x, y)] }$ असे समजा. $f$ ची पुढीलप्रमाणे व्याख्या करा:
\[
f(x) \simeq \umin{y}{R(x, y)}.
\]
मग $f$ आंशिक संगणनक्षम आहे आणि $S$ हे~$f$ चे परिभाषाक्षेत्र आहे.
\end{proof}


\end{document}
